%% kapa: exam=1 round=32 year=2021 period=1 form=A subject=civil_law subject_ko=민법 seq=28 number=28
%% source: Q-Net 감정평가사 기출문제 (cst00309) artlSeq=5208397 file=extracted/1008/5208397/A형/2021년도 제32회 감정평가사 1차 1교시 A형.pdf
%% conversion: kordoc 4.24.51 -> tools/gen1.py (auto); review=proofread-sonnet; flags=none
\begin{question}{28}
甲은 자신의 X토지를 乙에게 매도하였고, 乙은 X토지를 丙에게 전매하였다. 다음 설명으로 옳지 않은 것을 모두 고른 것은? (다툼이 있으면 판례에 따름)
\begin{framedtext}
ㄱ. 甲, 乙, 丙 사이에 중간생략등기에 관한 합의가 있다면, 甲의 乙에 대한 소유권이전등기의무는 소멸한다. \par
ㄴ. 乙의 甲에 대한 소유권이전등기청구권의 양도는 甲에 대한 통지만으로 대항력이 생긴다. \par
ㄷ. 甲, 乙, 丙 사이에 중간생략등기에 관한 합의가 없다면, 중간생략등기가 이루어져서 실체관계에 부합하더라도 그 등기는 무효이다. \par
ㄹ. 甲, 乙, 丙 사이에 중간생략등기에 관한 합의가 있은 후 甲·乙 간의 특약으로 매매대금을 인상한 경우, 甲은 인상된 매매대금의 미지급을 이유로 丙에 대한 소유권이전등기의무의 이행을 거절할 수 있다. \par
\end{framedtext}
\begin{choices}
  \item ㄱ, ㄴ
  \item ㄴ, ㄷ
  \item ㄷ, ㄹ
  \item ㄱ, ㄴ, ㄷ
  \item ㄴ, ㄷ, ㄹ
\end{choices}
\end{question}
\answer{4}
